% TOP_PROBLEM: {"id": 1350, "title": "GCH and Suslin Trees", "category_id": 10, "category": {"id": 10, "name": "set_theory", "display_name": "Set Theory", "description": "Foundations of mathematics, infinite sets, and cardinality.", "slug": "set-theory", "order_index": 10, "created_at": "2026-07-31T15:26:25.671Z"}, "status": "open"}
% FIRST_POSTED: 2026-10-01
% ATTEMPT_STATUS: unresolved
% Imported from: unsolvedmath_additional_500.tex; UnsolvedMath ID 1350
% !TeX program = lualatex
% Compile twice: lualatex 1350.tex
% Self-contained: no external bibliography, images, or \input files.
% Numeric section labels and visible numbers are original UnsolvedMath IDs.
\documentclass[11pt,a4paper]{article}
\usepackage[margin=24mm,headheight=14pt]{geometry}
\usepackage{fontspec}
\setmainfont{Latin Modern Roman}
\setsansfont{Latin Modern Sans}
\setmonofont{Latin Modern Mono}
\usepackage{amsmath,amssymb,mathtools}
\usepackage{unicode-math}
\setmathfont{Latin Modern Math}
\usepackage{microtype}
\usepackage{enumitem,xurl,needspace}
\usepackage[dvipsnames]{xcolor}
\usepackage{hyperref}
\hypersetup{unicode=true,colorlinks=true,linkcolor=MidnightBlue,urlcolor=MidnightBlue,
 pdftitle={UnsolvedMath Problem 1350},pdfauthor={Source-annotated compilation},bookmarksnumbered=true}
\usepackage{bookmark,tocloft,titlesec,fancyhdr}
\setlength{\cftsecnumwidth}{4.2em}
\cftsetpnumwidth{2.5em}
\cftsetrmarg{4em}
\titleformat{\section}{\Large\bfseries\raggedright}{\thesection}{0.7em}{}
\pagestyle{fancy}\fancyhf{}
\fancyhead[L]{\small\sffamily Additional UnsolvedMath statements}
\fancyhead[R]{\small Original catalogue IDs}
\fancyfoot[C]{\thepage}
\setlength{\parindent}{0pt}
\setlength{\parskip}{5pt plus 1pt}
\setlength{\emergencystretch}{3em}
\setcounter{tocdepth}{1}\setcounter{secnumdepth}{1}
\setlist[itemize]{leftmargin=3em,itemsep=4pt,topsep=4pt,parsep=0pt}
\allowdisplaybreaks
\newcommand{\N}{\mathbb{N}}
\newcommand{\Z}{\mathbb{Z}}
\newcommand{\Q}{\mathbb{Q}}
\newcommand{\R}{\mathbb{R}}
\newcommand{\C}{\mathbb{C}}
\newcommand{\F}{\mathbb{F}}
\newcommand{\A}{\mathbb{A}}
\newcommand{\PP}{\mathbb{P}}
\newcommand{\E}{\mathbb{E}}
\newcommand{\eps}{\varepsilon}
\newcommand{\dd}{\,\mathrm d}
\newcommand{\sref}[2]{\hyperlink{source:#1:#2}{[S#2]}}
\newif\ifshowcatalogue
\showcataloguetrue
\newcommand{\cataloguescope}[1]{\ifshowcatalogue
 \paragraph{Statement recorded in the supplied source.}
 #1\fi}
\begin{document}
% UnsolvedMath ID 1350; source catalogue code ST-004

\setcounter{section}{1349}

\section{Does GCH Force an Aleph-Two Suslin Tree?}\label{1350}

{\small\sffamily UnsolvedMath ID 1350 · Set Theory}

\subsection{Definitions and mathematical statement}

\paragraph{Shared notation.}Unless a section says otherwise, $\N=\{1,2,\ldots\}$,
$\Z$ is the ring of integers, $\Q,\R,\C$ are the rational, real, and complex
fields, $[N]=\{1,\ldots,\lfloor N\rfloor\}$, and $\log$ is the natural
logarithm. For real functions of a parameter tending to infinity,
$f\ll g$ and $f=O(g)$ mean $|f|\leq Cg$ eventually for some fixed $C$;
$f\gg g$ reverses this comparison; $f=o(g)$ means $f/g\to0$;
$f\sim g$ means $f/g\to1$; and $f\asymp g$ means both order bounds.
A subscript on $O$ or $\ll$ specifies allowed constant dependence.
The expression $n^{o(1)}$ in an upper bound means growth smaller than
$n^\epsilon$ for each fixed $\epsilon>0$. Local definitions govern any
exception, finite-field convention, or use of the same symbol for another
object. An asymptotic claim is not automatically an effective algorithm.



A $\kappa$-tree has height $\kappa$ and every level has size less than $\kappa$; the predecessors of each node are well-ordered. It is a $\kappa$-Suslin tree when it has neither a branch nor an antichain of size $\kappa$. The target is the implication, over ZFC, from GCH at all infinite cardinals to the existence of such a tree for $\kappa=\aleph_2$. \sref{1350}{1} \sref{1350}{2}

\paragraph{Statement recorded in the supplied source.}
Does the generalized continuum hypothesis imply the existence of an $\aleph_2$-Suslin tree? \sref{1350}{1}

\subsection{Short English statement}

Does the generalized continuum hypothesis guarantee a tree of height aleph-two with neither a very long branch nor a very large antichain?

\subsection{Research attempt: branch exclusion from splitting and the closure barrier}
\textbf{Status and context.} This supplies structural reductions and an obstruction to a naive construction, not an $\aleph_2$-Suslin tree under GCH. The inspected primary paper [R1] explicitly asks the same GCH question. Its additional square-based constructions are not treated as consequences of GCH alone.

\subsubsection{1. A cofinal branch creates an antichain in a splitting tree}
Consider a normal tree of regular height $\kappa$, with every node extending to every higher level and with two incompatible immediate successors at each node. If it has a cofinal branch, choose its node $b_\alpha$ at each level $\alpha<\kappa$. At the next level let $c_\alpha$ be an immediate successor different from the one followed by the branch. For $\alpha<\beta$, the node $c_\alpha$ is incompatible with $b_\beta$ and therefore with $c_\beta$. Thus $\{c_\alpha:\alpha<\kappa\}$ is an antichain of size $\kappa$.

Consequently, in this normal splitting setting, ruling out $\kappa$-sized antichains automatically rules out cofinal branches. A branch of cardinality $\kappa$ must be cofinal because each height below the initial ordinal $\kappa$ has size less than $\kappa$. This reduces the two prohibitions in the definition to the antichain prohibition for a suitably constructed tree; it does not produce that tree.

\subsubsection{2. The full binary tree breaks the width bound}
The full binary tree $2^{<\omega_2}$ splits everywhere, but its level $\omega_1$ has size $2^{\aleph_1}$. Under GCH this is $\aleph_2$, so it is not an $\aleph_2$-tree as defined in the question. Its many cofinal branches are another defect. Restricting attention to finite or countable levels misses the critical level $\omega_1$.

In fact, closure under all chains of length less than $\omega_2$ is incompatible with normal splitting and narrow levels, whether or not GCH holds. Suppose such a tree were given, with a root at level zero. Recursively embed every binary sequence of length at most $\omega_1$: at successor stages use the two distinct immediate successors, and at a limit stage take an upper bound of the previously chosen chain and then its predecessor at that limit level. The closure hypothesis supplies the upper bound; the definition of a normal tree supplies the predecessor. Distinct binary sequences stay incompatible after their first differing bit. At level $\omega_1$ this gives at least $2^{\aleph_1}\geq\aleph_2$ nodes, contradicting narrowness.

\subsubsection{3. The recursion has to prune rather than close indiscriminately}
The obstruction identifies a precise design constraint. A height-$\omega_2$ construction cannot preserve every branch through every limit level below $\omega_2$. At some limits it must discard enough candidate branches to keep width below $\omega_2$, while continuing every previously created node and anticipating potential large antichains. Simply taking the union of all earlier branches would reproduce the forbidden width in the preceding paragraph.

This does not rule out weaker closure. For example, closure only under countable chains does not require the upper bounds for chains of length $\omega_1$ used in the contradiction. Confusing those two closure strengths would falsely exclude standard higher-tree constructions. Likewise, the branch-to-antichain argument specifically requires splitting along the entire branch; a single chain has no such off-branch successors and is not a counterexample to that argument.

\subsubsection{4. Verification and exact residual problem}
Every cardinality step has been kept at its stated level: $2^{\aleph_1}\geq\aleph_2$ is Cantor's theorem in ZFC, while GCH merely identifies equality. The tree requirements are conditional hypotheses used to test a candidate construction; they have not been derived from the continuum function. No finite tree computation can certify height $\omega_2$ or the absence of an $\omega_2$-antichain. The missing result is a pruning or guessing principle provable from GCH alone that keeps all levels narrow and seals every prospective large antichain. The elementary reductions above do not supply that principle.

\subsection{Sources}

{\small

\begin{itemize}

\item[\hypertarget{source:1350:1}{S1}] ULAM.AI, \emph{UnsolvedMath}, supplied \texttt{problems.json}, record 1350 (ST-004), statement and background. The embedded literature-review date is 2026-08-17; that is the archive’s review date, not a fresh certification. Dataset landing page: \url{https://huggingface.co/datasets/ulamai/UnsolvedMath}.

\item[\hypertarget{source:1350:2}{S2}] Suslin tree overview. \url{https://en.wikipedia.org/wiki/Suslin_tree}. Bibliographic information imported from the supplied record.

\item[R1] A. Rinot, Higher Souslin trees and the GCH, revisited; primary author-hosted paper.
\url{https://www.assafrinot.com/files/rinot24.pdf}.
\end{itemize}

}

\label{end:1350}
\end{document}
